 

%\subsubsection{Data Augmentation.} \label{sec:data-augmentation}

%To allow the network learn the desired invariance and robustness properties, two strategies are applied to leverage the available annotated data (Fig. \ref{fig:data-aug}): 

%\begin{itemize}
 %   \item Flipping and rotating the image data by multiples of $\dfrac{\pi}{2}$.
    
  %  \item Rotating the input image by an angle sampled in $[-5^{\circ}, 5^{\circ}]$ following the uniform distribution. Subsequently, the transformed image is cropped in the created areas with no pixel data and resized to its original size, so there is no need to extrapolate over these areas.

%\end{itemize}

%\begin{figure}[H]
%\centering
%\includegraphics[width=4.8in]{Figures/data_augmentation.png}
%\caption{Data augmentation example used in segmentation. The same operations are applied to the corresponding mask.}
%\label{fig:data-aug}
%\end{figure}





%============================================================
%============================================================




\section{Experimental Results} \label{sec:results}

%In this section, we show the effectiveness of our segmentation approach through both quantitative and qualitative evaluations as well as showing a comparison with competing techniques, providing the results for every step in our algorithm.

In this section, we demonstrate the efficacy of our segmentation approach by conducting comprehensive quantitative and qualitative analysis. Additionally, we present a comparative assessment against competing techniques, offering detailed results for each step implemented within our algorithm.


%We now present the segmentation performance and computational cost of our proposed algorithm. To start with,  we examine the best hyperparameters reported in the BU-Net and random forest steps. The reliability of the algorithm is analyzed qualitatively and quantitatively, comparing each post-processing step. We also discuss whether our proposed algorithm can outperform the original U-Net. Finally, we cover the running time spent for each algorithm step.


%----------------------------------------------------------------------
\vspace{-0.5cm}
%\vspace{-0.3cm}
\subsection{Dataset} \label{sec:data} 
\vspace{-0.1cm}

We present a dataset that consists of high-resolution images taken by drone and ground-based equipment from distinct locations of the blade. The images provide a large variety of blade sizes, shapes, illumination conditions, as well as the background to be observed (Fig.~\ref{fig:introd}). Another background variation arises from on-shore and off-shore wind turbines. For this challenging dataset, we propose 1,712 images for training, 120 for validation and 200 for test. As a segmentation ground-truth is not provided, we manually annotated all the images. To make a fair analysis, the three sets are composed of different windfarms and inspection campaigns, ensuring their independence and that the test properly generalizes. In particular, we randomly selected 20 images of each windfarm to generate the validation and test sets. On the other hand, the training data was selected from a pool of different blade images, prioritizing the ones that are more challenging. These are the images in the root, the tip or the max-cord (where the blade section is widest). All color images and masks were min-max normalized. In addition to that, flipping and cropping strategies were applied to learn the desired invariance and robustness properties.



%The segmentation data consists in high-resolution images taken from distinct locations of the blade with their corresponding ground truth masks, which have been annotated manually. The dataset contains pictures taken by drone and ground-base equipment, which mainly diverge on its background. Other background variation arises from on-shore and off-shore wind turbines. The expected relation between background-blade pixels is 2:1, but in some cases the blade area could be very reduced. The data include 1712 training, 119 validation and 200 test images. The images from these three sets are from different windfarms and inspection campaigns, ensuring their independence and that the test properly generalizes.

 






%----------------------------------------------------------------
\vspace{-0.5cm}
%\vspace{-0.3cm}
\subsection{Implementation Details} \label{sec:training-seg}
\vspace{-0.1cm}

The original input images were resized to $1024 \times 1024$ pixels. We use Adam solver~\cite{adam} with an initial learning rate of $10^{-4}$ and a mini-batch of 1. A custom scheduler is employed to reduce the learning rate when there is no improvement on the validation loss after three epochs. Early stopping is performed. The model is trained using a NVIDIA GeForce RTX 3080 GPU. The weights for the focal loss in Eq.~\eqref{focal_loss} are set to $\gamma=2$ and $\alpha=0.25$, and the weight in Eq.~\eqref{total_loss} $\lambda=1$.



%----------------------------------------------------------------
\vspace{-0.5cm}
%\vspace{-0.3cm}
\subsection{Quantitative Evaluation}
\vspace{-0.1cm}

%The segmentation performance per image is analyzed in terms of the accuracy, recall and F1-score. 

%These metrics are defined pixel-wise as: $\textrm{accuracy} = \dfrac{TP + TN}{p}$, $\textrm{recall} =  \dfrac{TP}{TP+FN}$ and $\textrm{F1-score} = \dfrac{1}{\frac{1}{\text{precision}} + \frac{1}{\text{recall}}}$, where $TP$ indicates the number of True Positive pixels, $TN$ denotes the number of True Negative pixels, and $FN$ is the number of False Negative pixels. Moreover, precision is defined as $\frac{TP}{TP+FP}$ with $FP$ the number of False Positive pixels. 

The segmentation performance per image is analyzed in terms of accuracy, recall and F1-score. %The accuracy corresponds to the proportion of pixels of the predicted mask that match the ground truth mask, i.e., the ratio of correctly classified pixels. The recall is the ratio of blade pixels that are correctly classified by the algorithm. The F1-score is the harmonic mean of precision and recall, where the precision measures the accuracy of the blade predictions. The F1-score combines both precision and recall into a single metric, giving much more weight to low values.



%\begin{equation} \label{eq:accuracy}
 % \textrm{accuracy} = \dfrac{TP + TN}{p}, 
%\end{equation}


%\begin{equation} \label{eq:recall}
%  \textrm{recall} =  \dfrac{TP}{TP+FN}
%\end{equation}


%\begin{equation} \label{eq:f1}
%  \textrm{F1-score} = \dfrac{1}{\frac{1}{\text{precision}} + \frac{1}{\text{recall}}}
%\end{equation}

%where $TP = $ number of True Positive pixels, $TN = $ number of True Negative pixels, and $FN = $ number of False Negative pixels. Moreover, precision $= \frac{TP}{TP+FP}$ with $FP = $ number of False Positive pixels.




%the proportion of correctly classified blade pixels divided by the number of blade pixels in the ground truth mask.
%recall = true positive rate = true blade pixel rate


%The highest validation accuracy of the BU-Net is obtained using the original U-Net architecture, with a training objective $\mathcal L_{focal}(\alpha = 0.25, \gamma = 2) + \lambda  \cdot \mathcal L_{contiguity} $ where $\lambda = 1$, and a data augmentation strategy of small rotations followed by the corresponding zoom.


\textbf{Hyperparameter Search.} Our random-forest step (section~\ref{random_forest}) is optimized in terms of the number of split branches and the number of decision trees, obtaining the highest accuracy with five tree estimators that have a maximum of four split branches. Additionally, when the local pixel neighbors are included in the input, we enhance the random forest masks. However, increasing the number of input neighbors does not increase the segmentation accuracy, thus, we took as input the local pixel and its most immediate neighbors. 

\textbf{Test Set Results.} The performance over the test set is compared between each algorithm step. In particular, we report the accuracy, recall and F1-score in Table~\ref{tab:seg-accuracy}. 

\vspace{-1.5cm}
By only employing our BU-Net, we obtain a notable outcome of 92.28\% of accuracy. However, this result is not sufficiently good for blade assessments, because the recall of 85.73\% could compromise the detection of defects. For instance, if a large area of the blade is masked as background, it could cover the region where there are structural damages, which is a severe defect that cannot be overlooked.

By incorporating the first hole-filling step, we substantially improve performance. Namely, the recall increases to 92.08\%, which represents the highest gain of all the post-processing steps. This huge difference is due to the BU-Net being able to successfully capture the blade borders, but not the inner region of the blade. Further algorithm steps are necessary to yield further improvement, reaching in the end a novel performance of 97.39\% of accuracy and 93.35\% of recall. On balance, and as it can be seen in Table~\ref{tab:seg-accuracy}, every step is key to sort out the problem properly, obtaining a final accurate solution that cannot be achieved directly by our BU-Net. 

Fig.~\ref{fig:boxplot} analyzes the performance dispersion of each step. We observe the skewness of the BU-Net performance: rather spread along the tail with lower performance, containing many instances with poorer accuracy and recall compared with their median. As we include the different steps in our algorithm, that dispersion is reduced notably. Therefore, despite not substantially improving the mean performance, the subsequent steps are the cornerstone to ensure robustness.

\begin{table}[t!]
\centering
\caption{\textbf{Test results after each segmentation step.} }
    \vspace{-0cm}
    \label{tab:seg-accuracy}
%\begin{tabular}{p{1.95cm}p{1.0cm}p{1.2cm}p{1.6cm}p{1.1cm}}
\resizebox{8.5 cm}{!} {
\begin{tabular}{ccccc}
\toprule
     \multicolumn{1}{c}{Performance} & \multicolumn{1}{c}{BU-Net} &   \multicolumn{1}{c}{Fill holes}  &  \multicolumn{1}{c}{Random forest}  &  \multicolumn{1}{c}{Fill holes 2}\\
    \multicolumn{1}{c}{metric} &  \multicolumn{1}{c}{${\hat{\mathbf{S}}}^{BU}$} &   \multicolumn{1}{c}{$\hat{\mathbf{S}}^{H1}$} &  \multicolumn{1}{c}{$\hat{\mathbf{S}}^{RF}$} &  \multicolumn{1}{c}{$\hat{\mathbf{S}}^{H2}$} \\ \midrule
   Accuracy (\%) &  92.28 & 96.11 & 97.11 & 97.39\\
   Recall (\%) &  85.73 & 92.08 & 92.37 & 93.35 \\
   F1-score (\%) & 90.41 & 94.53 & 95.12 & 95.73 \\
   mIoU (\%) & 87.71 & 93.57 & 95.07 &  95.53 \\
\bottomrule
\end{tabular}}
\vspace{-0.25cm}
\end{table}

\begin{figure}[t!] %[htb]
\begin{minipage}[b]{1.0\linewidth}
  \centering
  \resizebox{8 cm}{!} {
  \centerline{\includegraphics[width=8.5cm]{Figures/Results/recall_boxplot.png}}}
  \vspace{-0.35cm}
\end{minipage}
\caption{\textbf{Box plot test results after each segmentation step.} }
\label{fig:boxplot}
\vspace{-0.4cm}
\end{figure}

% \begin{table}[t!]
% \centering
% \caption{\textbf{Test results after applying each step in our segmentation algorithm.} }
%     \vspace{-0cm}
%     \label{tab:seg-accuracy}
% \begin{tabular}{p{1.95cm}p{1.0cm}p{1.2cm}p{1.6cm}p{1.1cm}}
% %\begin{tabular}{ccccc}
% \toprule
%      \multicolumn{1}{p{1.95cm}}{Performance} & \multicolumn{1}{p{1.0cm}}{BU-Net} &   \multicolumn{1}{p{1.2cm}}{Fill holes}  &  \multicolumn{1}{p{1.6cm}}{Random forest}  &  \multicolumn{1}{p{1.1cm}}{Fill holes 2}\\
%     \multicolumn{1}{p{1.95cm}}{metric} &  \multicolumn{1}{p{1.0cm}}{${\hat{\mathbf{S}}}^{BU}$} &   \multicolumn{1}{p{1.2cm}}{$\hat{\mathbf{S}}^{H1}$} &  \multicolumn{1}{p{1.6cm}}{$\hat{\mathbf{S}}^{RF}$} &  \multicolumn{1}{p{1.1cm}}{$\hat{\mathbf{S}}^{H2}$} \\ \midrule
%    Accuracy (\%) &  92.28 & 96.11 & 97.11 & 97.39\\
%    Recall (\%) &  85.73 & 92.08 & 92.37 & 93.35 \\
%    F1-score (\%) & 90.41 & 94.53 & 95.12 & 95.73 \\
% \bottomrule
% \end{tabular}
% \vspace{-0.3cm}
% \end{table}

% \begin{table}[t!]
% \centering
% \caption{\textbf{Test results after applying each step in our segmentation algorithm.} }
%     \vspace{-0cm}
%     \label{tab:seg-accuracy}
% %\begin{tabular}{p{1.95cm}p{1.0cm}p{1.2cm}p{1.6cm}p{1.1cm}}
% \begin{tabular}{ccccc}
% \toprule
%      Performance & BU-Net &  Fill holes  &  Random forest  &  Fill holes 2\\
%     metric &  ${\hat{\mathbf{S}}}^{BU}$ &   $\hat{\mathbf{S}}^{H1}$ &  $\hat{\mathbf{S}}^{RF}$ & $\hat{\mathbf{S}}^{H2}$ \\ \midrule
%    Accuracy (\%) &  92.28 & 96.11 & 97.11 & 97.39\\
%    Recall (\%) &  85.73 & 92.08 & 92.37 & 93.35 \\
%    F1-score (\%) & 90.41 & 94.53 & 95.12 & 95.73 \\
% \bottomrule
% \end{tabular}
% \vspace{-0.3cm}
% \end{table}

%submitted
%variants 
%currently the 
%top performing method on the PASCAL
%report
%conducted 
%attaining
%focus
%yields
%significantly outperforming
%further improves performance to 63.1 percent. 
%observe any improvement when adopting either
%Employing the dense
%substantially outperforms the concurrent work 

% In Fig.~\ref{fig:seg-recall}, we analyze the performance dispersion of each step in our algorithm. We observe the skewness of the BU-Net performance: it is rather spread along the tail with lower performance, containing many instances with poorer accuracy and recall compared with their median. Note that once again the recall performance presents the worst results, with a lower bound extreme at 56.63\%. As we include the different steps in our algorithm, that dispersion is reduced notably. Therefore, despite not substantially improving the mean performance, the subsequent steps are the cornerstone to ensure the robustness of our algorithm, obtaining a final lower bound extreme of 88.57\% of recall.

% %Note that the whiskers are not actually at 1.5 x IQR from the box, they are the min and max data values within that range. 

% \begin{figure}[t!]
% \centering
% \includegraphics[width=3.45in]{Figures/Results/recall_boxplot.png}
%     \vspace{-0.3cm}
%     \caption{\textbf{Box plot test results after applying each segmentation step.} Performance dispersion of every algorithm step.}
%     \label{fig:seg-recall}
%     \vspace{-0.1cm}
% \end{figure}

\begin{table}[t!]
\centering
\caption{\textbf{Quantitative comparison on blade segmentation.} Relative metrics are with respect to the worst estimation.} 
\vspace{0cm}
\label{tab:seg-unet-compare}
\resizebox{8.5 cm}{!} {
\begin{tabular}{ccccccc}
\toprule
     \multicolumn{1}{c}{Method} & \multicolumn{1}{c}{Accuracy}  & \multicolumn{1}{c}{Recall} & \multicolumn{1}{c}{F1-score} & \multicolumn{1}{c}{mIoU} & Relative & Relative\\
    \multicolumn{1}{c}{} & {\%} & {\%}  & {\%}  & {\%} & accuracy & recall \\
    \midrule
   U-Net~\cite{unet} & 86.24 & 68.93 & 77.95 & 79.94 & 1 & 1\\ 
   DeepLabv3+~\cite{deeplabv3+}  & 94.14 & 87.38 & 89.03 & 90.74 & 1.09 & 1.27 \\ 
   SW~\cite{sw} & 93.48 & 91.71 & 91.37 & 89.71 & 1.08 & 1.33 \\ 
   ResNeSt~\cite{resnest} & 94.23 & 91.47 & 92.77 & 91.20 & 1.09 & 1.33 \\ 
   U-NetFormer~\cite{unetformer} & 96.20 & \textbf{93.51} & 94.42 & 93.65 & 1.12 & \textbf{1.36} \\ 
   Ours $\hat{\mathbf{S}}$ & \textbf{97.39} & 93.35 & \textbf{95.73} & \textbf{95.53} & \textbf{1.13} & 1.35 \\ 
\bottomrule
\end{tabular}}
\vspace{-0.5cm}
\end{table}



\textbf{Comparative Results.} As the BU-Net employs the U-Net architecture~\cite{unet}, we compared our proposed algorithm with the original one and state-of-the-art methods (Table~\ref{tab:seg-unet-compare}). We outperform those methods, thanks to overcoming the blade-background imbalances, and escaping from the tendency of predicting the most common region. 

%-----------------------------------------------------------------
\vspace{-0.3cm}
\subsection{Qualitative Results}
\vspace{-0.1cm}

We visualize the segmentation results of each algorithm step in Fig.~\ref{fig:seg-visual}. In general, our BU-Net can capture the blade edges and, therefore, its overall shape. However, it struggles with the inner area of the blade, which is addressed in the first hole filling step. Note that just using the standard hole filling approach would not improve the BU-Net outputs. Then, the random forest is responsible for fixing the misplaced pixels, acting as a regularizer of the pixel color intensity. By training a simple model, we ensure there is no overfitting and that it just learns the main patterns. Hence, if the BU-Net can in general classify which color intensities belong to the blade region, then the random forest would learn those patterns and resolve the particular images or pixels that failed, obtaining robustness in the solution. The second hole filling step is applied to solve the shortcomings the random forest can introduce. They are generally caused by dark regions that highly differ with the typical white color of the blade --a typical example is illustrated in Fig.~\ref{fig:seg-visual}, first row/right side--, where the random forest misplaces the vortex generators as background.

%the main patterns from all the images of the same blade surface, acting as a denoising step. In short, the random forest role is to tackle the specific errors of a single image. In the end, due to the simplicity of the random forest as a model

\begin{figure}[t!]
\resizebox{8.5 cm}{!} {
\hspace{-0.1cm}\begin{tabular}{@{}c|c@{}}
\includegraphics[width=0.4\textwidth]{Figures/Results/seg_visual_notruth_1.png}&
\includegraphics[width=0.4\textwidth]{Figures/Results/seg_visual_notruth_7.png}\\
%\includegraphics[width=0.4\textwidth]{Figures/Results/seg_visual_notruth_2.png}\\
\includegraphics[width=0.4\textwidth]{Figures/Results/seg_visual_notruth_3.png}&
\includegraphics[width=0.4\textwidth]{Figures/Results/seg_visual_notruth_4.png}\\
\includegraphics[width=0.4\textwidth]{Figures/Results/seg_visual_notruth_5.png}&
\includegraphics[width=0.4\textwidth]{Figures/Results/seg_visual_notruth_6.png}\\
%\includegraphics[width=0.4\textwidth]{Figures/Results/seg_visual_notruth_7.png}&
%\includegraphics[width=0.4\textwidth]{Figures/Results/seg_visual_notruth_8.png}
\end{tabular}}
\vspace{-0.2cm}
\caption{\textbf{Qualitative evaluation on test images after applying each step of our proposed segmentation algorithm.} On both sides, the same information is displayed. \textbf{First column:} Input color image ($\mathbf{I}$). \textbf{From second to fifth column:} BU-Net (${\hat{\mathbf{S}}}^{BU}$), first hole filling ($\hat{\mathbf{S}}^{H1}$), random forest ($\hat{\mathbf{S}}^{RF}$) and second hole filling ($\hat{\mathbf{S}}^{H2}\equiv\hat{\mathbf{S}}$) estimations, respectively.}
\label{fig:seg-visual}
\vspace{-0.35cm}
\end{figure}


In Fig.~\ref{fig:seg-fail}, we depict the most common cases where our segmentation algorithm could fail. To start with, the BU-Net might occasionally struggle with images of the wind turbine hub, because our dataset has a few instances of the hub. This is the case for the images in the first row. Distinct sources of contamination can also cause the BU-Net to mislead part of the blade with the background. This is the case for the image in the second row and right side. As the blade image borders are not fully captured and these images may contain a broad range of color intensities in the blade region, the hole filling and random forest steps are not sufficient to recover the blade region that has not been identified by the BU-Net.

In addition, an image could include wind turbines at its background or its own tower (see images in the second row). Ideally, we would like to predict the tower as blade, since it belongs to the turbine itself, but not the turbines in the landscape. This challenging problem is part of our future work. 


%Ideally, we would like to predict the tower as mask, since it belongs to the turbine itself. Turbines in the background of the image, however, should not be part of the blade. This challenging problem is part of our future work.



\begin{figure}[t!]
\resizebox{8.5 cm}{!} {
\hspace{-0.1cm}
\begin{tabular}{@{}c|c@{}}
\includegraphics[width=0.4\textwidth]{Figures/Results/seg_visual_fail_notruth_1.png}&
\includegraphics[width=0.4\textwidth]{Figures/Results/seg_visual_fail_notruth_2.png}\\
\includegraphics[width=0.4\textwidth]{Figures/Results/seg_visual_fail_notruth_3.png}&
\includegraphics[width=0.4\textwidth]{Figures/Results/seg_visual_fail_notruth_4.png}
\end{tabular}}
\vspace{-0.1cm}
\caption{\textbf{Failure cases on test images after applying each step of our segmentation algorithm.} See caption of Fig.~\ref{fig:seg-visual}.}
\label{fig:seg-fail}
\vspace{-0.0cm}
\end{figure}
%In both cases, the same information is displayed. \textbf{First column:} Input color image ($\mathbf{I}$). \textbf{From second to fifth column:} BU-Net (${\hat{\mathbf{S}}}^{BU}$), first hole filling ($\hat{\mathbf{S}}^{H1}$), random forest ($\hat{\mathbf{S}}^{RF}$) and second hole filling ($\hat{\mathbf{S}}^{H2}$) estimations, respectively.

%-----------------------------------------------------------
\vspace{-0.2cm}
\subsection{Windfarm Dissimilarity}
\vspace{-0.1cm}

%As mentioned in Section~\ref{sec:data}, as each dataset includes images from distinct windfarms, we examine in Fig.~\ref{fig:seg-inspections} how the performance varies depending on this factor.

%In each step of the algorithm, we observe significant dissimilar performance depending on the windfarm. For instance, in the first step when the BU-Net performance is reported, Windfarm \#1 performance is quite poorly compared with its piers, achieving approximately an accuracy of 83\%. In addition to that, each step of the algorithm can have drastically different effects between two windfarms. The extreme example comes again with Windfarm \#1, which after the first fill holing, it goes from the lowest accuracy to the highest performance. Conversely, for other cases such as Windfarm \#6, the first fill holing step does not signify any gain. A typical example from Windfarm \#1 that illustrates this phenomena is shown in Fig.~\ref{fig:seg-visual} (second row/ right side), where the BU-Net only captures the blade borders and then the hole filling is able to fill accurately the inner region of the blade. Another clear example comes with the random forest step, where while it significantly increases the accuracy of the Windfarm \#2, \#3 and \#7, it negatively affects the performance of Windfarm \#4 and \#5. An illustrative example of its improvement is shown in Fig.~\ref{fig:seg-visual} (first row/ left side).

Fig.~\ref{fig:seg-inspections} proves the generability of our approach, demonstrating that images from all the distinct windfarms have a high performance. Notice that this is not the case for the BU-Net without the post-processing steps, as it has really low performance for Windfarm \#1. Hence, it is the post-processing steps as a whole which ensures the generability. For instance, hole filling is crucial for Windfarm \#1, while random forest is crucial for Windfarm \#2 and \#3.
\vspace{-0.2cm}
%aquí resultados en el apéndice?

\begin{figure}[t!]
\centering
\includegraphics[width=3.1in]{Figures/Results/inspections.png}
    \vspace{-0.3cm}
    \caption{\textbf{Accuracy results over the test set for each windfarm.} It is included in each step in our method. For readability, the plot does not include every windfarm of the test set.}
    \label{fig:seg-inspections}
    \vspace{-0.3cm}
\end{figure}

%20210427-EUROWIND-DOSTRUP/Prise3200399T4-B-WW-pZR_001.jpg W7
%20210428-EUROWIND-DOSTRUP/Prise3200398T3-B-WW-tC3_010.jpg W2
%20211117-MDU-Diamond Hill/PriseTurbine07 Re-G-WW-DAw_010.jpg 
%Pioneer-Prairie-II/123 2S (119).JPG
%edf-turbine-2/B_WW_20155131003_backup.jpg
%edf-turbine-2/G_TE_20155131009_backup.jpg
%loegtved/a8ffa467-fd6e-49c8-9f1e-0d7f700bd72e_0044.jpg
%loegtved/a8ffa467-fd6e-49c8-9f1e-0d7f700bd72e_0078.jpg

%20210427-EUROWIND-DOSTRUP/Prise3200400T5-B-LE-XJr_001.jpg W7
%20211116-MDU-Diamond Hill/PriseTurbine12-R-LW-gFe_001.jpg
%Meridian-Way-II-207/207 2L (057).JPG W6
%upteko-ewii-marstrup/6_turbineSix_C_SS_0014.JPG







